\section{Introduction}
\label{sec:intro}

When diffusion generative models were first developed \cite{SohlDickstein2015,Song2019,Ho2020}, the core idea was supposed to be \mbox{\textit{denoising}}, \ie, \textit{predicting a clean image} from its corrupted version. 
However, two significant milestones in the evolution of diffusion models turned out to deviate from the goal of directly predicting clean images. 
First, predicting the noise itself (known as ``$\e$-prediction'') \cite{Ho2020} made a pivotal difference in generation quality and largely popularized these models. Later, diffusion models were connected to flow-based methods \cite{lipman2022flow,liu2022flow,albergo2023building} through predicting the flow velocity (``\mbox{$\v$-prediction}'' \cite{salimans2022progressive}), a quantity that combines clean data and noise.
Today, diffusion models in practice commonly predict noise or a noised quantity (\eg, velocity).

% ##################################################
\begin{figure}[t]
\centering
\includegraphics[width=0.8\linewidth]{figures/teaser-update}
\vspace{-.5em}
\caption{
\textbf{The Manifold Assumption} \cite{Chapelle2006SemiSupervised}
hypothesizes that natural images lie on a low-dimensional manifold within the high-dimensional pixel space.
While a clean image $\x$ can be modeled as on-manifold, the noise $\e$ or flow velocity $\v$ (\eg, $\v = \x - \e$) is inherently off-manifold. Training a neural network to 
predict a clean image (\ie, $\x$-prediction) is \textit{fundamentally different} from training it to predict noise or a noised quantity (\ie, $\e$/$\v$-prediction).
}
\label{fig:teaser}
\vspace{-1em}
\end{figure}
% ##################################################

Extensive studies \cite{salimans2022progressive,Karras2022edm,hoogeboom2023simple,esser2024scaling} have shown that predicting the clean image (``\mbox{$\x$-prediction}" \cite{salimans2022progressive}) is closely related to $\e$- and $\v$-prediction, provided that the \textit{weighting} of the prediction loss is reformulated accordingly (detailed in \cref{sec:analysis}).
Owing to this relation, less attention has been paid to what the network should directly predict, implicitly assuming that the network is capable of performing the assigned task.

However, the roles of clean images and a noised quantity (including the noise itself) are \textit{not} equal.
In machine learning, it has long been hypothesized \cite{Chapelle2006SemiSupervised,carlsson2009topology} that
\textit{``(high-dimensional) data lie (roughly) on a low-dimensional manifold''} (\cite[p.6]{Chapelle2006SemiSupervised}).
Under this manifold assumption, while clean data can be modeled as lying on a low-dimensional manifold, a noised quantity is inherently distributed across the full high-dimensional space \cite{vincent2010stacked} (see \cref{fig:teaser}). \textit{Predicting clean data is fundamentally different from predicting noise or a noised quantity.}

Consider a scenario where a low-dimensional manifold is embedded in a high-dimensional observation space.
Predicting noise in this high-dimensional space requires \textit{high capacity}: the network needs to preserve all information about the noise.
In contrast, a \textit{limited-capacity} network can still predict the clean data, as it only needs to retain the low-dimensional information while filtering out the noise.

When a low-dimensional space (\eg, image latent \cite{Rombach2022}) is used, 
the difficulty of predicting noise is alleviated, yet at the same time is \textit{hidden} rather than solved.
When it comes to pixel or other high-dimensional spaces, existing diffusion models can still struggle to address the \textit{curse of dimensionality} \cite{Chapelle2006SemiSupervised}.
The heavy reliance on a pre-trained latent space  prevents diffusion models from being self-contained.

In pursuit of a self-contained principle, there has been strong focus on advancing diffusion modeling in the pixel space \cite{chen2023importance,hoogeboom2023simple,hoogeboom2024simpler,chen2025pixelflow,wang2025pixnerd}. 
In general, these methods explicitly or implicitly avoid the information bottleneck in the networks, \eg, by using dense convolutions or smaller patches, increasing channels, or adding long skip connections. 
We suggest that these designs may stem from the demand to predict high-dimensional noised quantities.

In this paper, we return to first principles and let the neural network \textit{directly} predict the clean image. By doing so, we show that a \textit{plain} Vision Transformer (ViT) \cite{Dosovitskiy2021}, operating on large image patches consisting of raw pixels, can  be effective for diffusion modeling.
Our approach is \textit{self-contained} and does not rely on any pre-training or auxiliary loss --- \textit{no latent tokenizer \cite{Rombach2022}, no adversarial loss \cite{goodfellow2014generative,Rombach2022}, no perceptual loss \cite{zhang2018unreasonable,Rombach2022} (thus no pre-trained classifier \cite{simonyan2014very}), and no representation alignment \cite{repa} (thus no self-supervised pre-training \cite{oquab2023dinov2})}. 
Conceptually, our model is nothing more than ``Just image Transformers'', or \textbf{JiT}, as we call it, applied to diffusion.

We conduct experiments on the ImageNet dataset \cite{deng2009imagenet} at resolutions of 256 and 512, using JiT models with patch sizes 16 and 32 respectively. 
Even though the patches are very high-dimensional (hundreds or thousands), our models using $\x$-prediction can easily produce strong results, where $\e$- and $\v$-prediction fail catastrophically.  Further analysis shows that it is \mbox{unnecessary} for the network width to match or exceed the patch dimension; in fact, and surprisingly, a \textit{bottleneck} design can even be beneficial, echoing observations from classical manifold learning.

Our effort marks a step toward a self-contained ``Diffusion\,+\,Transformer'' philosophy \cite{Peebles2023} on native data. Beyond computer vision, this philosophy is highly desirable in other domains involving natural data (\eg, proteins, molecules, or weather), where a tokenizer can be difficult to design. By minimizing domain-specific designs, we hope that the general ``Diffusion\,+\,Transformer'' paradigm originated from computer vision will find broader applicability.